\begin{diagram}[
  type=flow,
  label={fig:sec06-quality-loop},
  caption={A failing test does not pass silently: Test branches to Publish only once checks pass, while a failure follows a dashed exception path to Alert and then to an explicit block, warn, or quarantine decision that revises the Contract for the next interval.},
  source={Author's synthesis.},
  description={A five-node flow with one branch point. Contract leads to Test. Test is the branch point: on pass, a solid edge leads to Publish, a terminal state; on fail, a dashed exception edge labeled fail leads to Alert. Alert leads by a dashed edge labeled quarantine to a Remediate or Quarantine decision, where the action is to block publication, warn the owner, or quarantine the affected data rather than a silent pass. A dashed edge then loops back from Remediate or Quarantine to Contract, closing the recurring cycle with no fixed start or end.}
]
% Main row is the solid data path (Contract -> Test -> Publish on pass).
% The lower row is the dashed exception path (fail -> quarantine -> loop back),
% positioned directly under Contract so the return edge stays clear of Test.
\stepat[text width=24mm,align=center]{contract}{0}{0}{Contract}
\stepat[text width=24mm,align=center]{test}{3.15}{0}{Test}
\stepat[rk terminal,text width=24mm,align=center]{publish}{6.30}{0}{Publish}
\stepat[text width=24mm,align=center]{alert}{3.15}{-2.2}{Alert}
\stepat[text width=24mm,align=center]{remediate}{0}{-2.2}{Remediate or\\Quarantine}
\flowedge{contract}{test}
\flowedge{test}{publish}
\RKEdge[fail]{optional}{test}{alert}
\RKEdge[quarantine]{optional}{alert}{remediate}
\RKEdge{optional}{remediate}{contract}
\end{diagram}
